\subsubsection{1D Distributed Area--Flow Models}
\label{subsubsec:1d-models}

The selected one-dimensional area--flow model is specified in
Section~\ref{subsec:methodology-problem-definition}. There the retained fields
$A(z,t)$, $Q(z,t)$, $\bar u(z,t)=Q/A$, and $p_{\mathrm g}(z,t)$, the
elastic-wall area--flow balance law, the stenosis reference geometry, the
conservative flux/source notation, and the single-vessel boundary data are
recorded in
Definitions~\ref{def:section-averaged-1d-fields}--\ref{def:1d-single-vessel-boundary-data}
\cite{FormaggiaEtAl2003OneDimensionalBloodFlow,KopplHelmig2023DimensionReducedModeling}.
This subsection uses that notation only to identify the hierarchy role of the
1D tier.

In the hierarchy, the key point is that the 1D state is distributed but reduced.
It retains axial area, flow, mean velocity, and gauge pressure; it does not
retain the full three-dimensional velocity field, pressure field, secondary
flow, separation, recirculation, or wall traction. Reduced shear and viscosity
quantities therefore inherit the declared velocity-profile, wall-friction, and
rheology closures. Tube-flow and hemorheology data caution that diameter,
hematocrit, and shear-rate dependence cannot be promoted to resolved local
physics without a declared closure map
\cite{PriesEtAl1992BloodViscosityTubeFlow,GaldiEtAl2008HemodynamicalFlows}.
Reduced shear proxies are not traction-based WSS unless a separate observation
map declares how a reduced output is compared with a resolved wall-stress
quantity.

The closure contract remains part of the model definition. A 1D case must
declare the momentum-flux/profile convention, friction law, rheology descriptor
$\mathcal R$, wall descriptor $\mathcal W$, coupling descriptor $\mathcal C$,
boundary data, numerical realization, and output operators. Stenosis enters
through the reference geometry $R_0(z)$, $A_0(z)$, and $s(z)$, but resolved
throat flow does not appear automatically; variable-radius corrections such as
the Canic-style terms are additional closure choices
\cite{CanicEtAl2024Extended1DStenosis}.
The current implementation records this axis explicitly through the
\texttt{canic-extended-1d} forward-model descriptor and the historical
\texttt{classical-1d-no-slip} manifest token for the classical
parabolic-profile baseline. Both use the selected
\texttt{canic-koiter-thin-membrane} pressure--area wall law; the former adds
the Canic variable-radius corrections, while the latter corresponds to the
parabolic-profile baseline in
Definition~\ref{def:classical-no-slip-1d-baseline}.

The 1D model is the selected reduced forward map because it retains distributed
pressure--flow structure at low computational cost, while its outputs remain
reduced quantities with declared observation maps. In this report it is the
selected single-vessel generator for the numerical study, not a
clinical-validity claim, an independent-solver parity claim, or a
resolved-3D-accuracy claim.
